\documentclass[10pt,letterpaper]{article}

\usepackage[margin=0.58in]{geometry}
\usepackage[T1]{fontenc}
\usepackage[utf8]{inputenc}
\usepackage[default,type1]{sourcesanspro}
\usepackage[final]{microtype}
\usepackage[explicit]{titlesec}
\usepackage{enumitem}
\usepackage{tabularx}
\usepackage{array}
\usepackage{ragged2e}
\usepackage{iftex}
\usepackage[hidelinks]{hyperref}

\ifPDFTeX
  \input{glyphtounicode}
  \pdfgentounicode=1
\fi

\hypersetup{
  pdftitle={Connor Savugot Resume},
  pdfauthor={Connor Savugot},
  pdfsubject={Embedded Firmware Engineering},
  pdfkeywords={embedded firmware, C, C++, Embedded Linux, BMS, CAN, I2C, SPI, board bring-up}
}

\pagestyle{empty}
\setcounter{secnumdepth}{0}
\setlength{\parindent}{0pt}
\setlength{\parskip}{0pt}
\hfuzz=0.4pt

\titleformat{\section}
  {\fontsize{11.5}{12.2}\selectfont\bfseries\uppercase}
  {}{0pt}{#1\hspace{0.12cm}\titlerule[0.65pt]}
\titlespacing{\section}{0pt}{0.19cm}{0.08cm}

\newlist{resumebullets}{itemize}{1}
\setlist[resumebullets]{
  label=\textbullet,
  leftmargin=0.43cm,
  topsep=0.02cm,
  itemsep=0.018cm,
  parsep=0pt,
  partopsep=0pt
}

\newcommand{\companyheader}[1]{%
  {\fontsize{10.7}{11.2}\selectfont\bfseries #1}\par\vspace{-0.02cm}
}

\newcommand{\resumeentry}[2]{%
  \begin{tabularx}{\textwidth}{@{}X>{\raggedleft\arraybackslash}p{2.50in}@{}}
    \textbf{#1} & #2
  \end{tabularx}\vspace{-0.06cm}
}

\newcommand{\roleentry}[2]{%
  \begin{tabularx}{\textwidth}{@{}X>{\raggedleft\arraybackslash}p{2.50in}@{}}
    \textit{#1} & #2
  \end{tabularx}\vspace{-0.06cm}
}

\begin{document}
\fontsize{10}{11.15}\selectfont
\setlength{\RaggedRightRightskip}{0pt plus 5em}
\RaggedRight
\emergencystretch=1em
\hyphenpenalty=10000
\exhyphenpenalty=10000

\begin{center}
  {\fontsize{21.5}{22.5}\selectfont\bfseries Connor Savugot}\\[2pt]
  {\bfseries B.S. Computer Engineering | North Carolina State University | May 2026}\\[2pt]
  {\fontsize{9.7}{10.5}\selectfont
    \href{mailto:consav04@gmail.com}{consav04@gmail.com}
    \enspace|\enspace \href{tel:+18285419487}{+1-828-541-9487}
    \enspace|\enspace Murphy, NC
    \enspace|\enspace \href{https://linkedin.com/in/connorsavugot}{linkedin.com/in/connorsavugot}
    \enspace|\enspace \href{https://github.com/Clem085}{github.com/Clem085}}
\end{center}
\vspace{-0.16cm}

\section{Technical Profile}
Embedded firmware engineer delivering programmable power-conversion control, BMS and IC integration, Embedded Linux services, board bring-up, and protocol-level debugging from architecture and datasheet evaluation through custom-PCB validation.

\section{Technical Skills}
\textbf{Firmware \& Systems:} C, C++, Python, Bash; FreeRTOS, Embedded Linux, Yocto/Arago, systemd\\
\textbf{Hardware \& Interfaces:} TI TMS320, TI AM62, MSP430FR2355; I2C, SPI, UART/RS232, CAN/CAN FD, ADC, PWM, GPIO, IPMI/IPMB\\
\textbf{Development \& Test:} Git/GitLab, Code Composer Studio, KiCad, SocketCAN/PCAN, JTAG, oscilloscopes, logic/protocol analyzers, through-hole/SMT soldering

\section{Engineering Experience}
\companyheader{AEGIS Power Systems}
\roleentry{Embedded Firmware Engineer}{Present}
\begin{resumebullets}
  \item Implemented and debugged CAN/CAN FD, I2C, SPI, and UART interfaces; isolated firmware, timing, wiring, configuration, and physical-layer faults with register-level tests, scopes, analyzers, known-good nodes, and controlled substitutions.
  \item Supported TI AM62 Yocto/Arago BSPs and device trees; wrote systemd and external-watchdog services, built Python/Bash diagnostics, and validated IPMI/IPMB sensor and FRU access.
\end{resumebullets}

\roleentry{Engineering Intern}{May--Aug 2024 \& May--Aug 2025}
\begin{resumebullets}
  \item Developed C firmware for custom bidirectional buck--boost power supplies across transformer-coupled stages using 12 PWM generator pairs (24 physical A/B outputs); independently configured each A output's switching period/frequency, duty cycle, and relative phase, then set its paired B output to match or exactly invert A.
  \item Built a restricted SSH/chroot command interface and Python/Bash tools for approved remote control and inspection of TI Yocto/Arago embedded power hardware.
  \item Developed packetized, checksum-validated TMS320 flashing utilities, reverse-engineered RS232 command/checksum transactions from limited documentation, and supported FreeRTOS, serial-interface, schematic, and burn-in test debugging.
\end{resumebullets}

\resumeentry{TEAM Industries | Engineering Intern}{Jun--Aug 2021 \& Jun--Aug 2022}
\begin{resumebullets}
  \item Built an Excel/VBA tolerance checker for out-of-spec manufactured parts and organized a searchable 17,000+ record catalog linking tooling, parts, vendors, operations, locations, and revision data.
\end{resumebullets}

\section{Systems Design}
\resumeentry{EV Active Sensor Adapter | MSP430FR2355}{NC State University}
\begin{resumebullets}
  \item Defined architecture, requirements, interface specifications, and validation plans for an embedded system spanning sensing, BMS/power-management hardware, custom-PCB interfaces, communications, and telemetry.
  \item Evaluated BMS/support ICs and prototyped interrupt-driven SPI, ADC, UART, timer, GPIO, and LCD interfaces while coordinating integration through trade studies, risk tracking, design reviews, and subsystem test plans.
  \item Evaluated BMS and power-management ICs against electrical, protocol, packaging, and sourcing constraints, then carried selected devices through breakout-board testing, driver/peripheral bring-up, custom-PCB integration, and validation.
\end{resumebullets}

\section{Leadership \& Teaching}
\resumeentry{Embedded Systems (ECE306) Teaching Assistant | NC State University}{}
\begin{resumebullets}
  \item Delivered supplemental lectures and laboratory support for C/MSP430 firmware, ADC, timers, interrupts, GPIO, LCDs, and custom-PCB bring-up; taught systematic diagnosis of code, wiring, soldering, and toolchain faults.
\end{resumebullets}

\resumeentry{Lead Instructor | IEEE Open Project Space 2}{Two Semesters}
\begin{resumebullets}
  \item Led a ten-week, project-based electronics course from scope through final demonstration; taught through-hole/SMT soldering with custom IEEE PCBs and supported student development using PlatformIO, GitHub, KiCad, and BOM workflows.
\end{resumebullets}

\end{document}
